\begin{enumerate}
\item Nắm chắc các khái niệm phép thử, không gian mẫu, biến cố (là tập con của $\Omega$), biến cố chắc chắn, biến cố không thể và biến cố đối.
\item Xét mỗi khẳng định: đối chiếu với định nghĩa, hoặc lấy một phép thử cụ thể (gieo xúc xắc, tung đồng xu) để kiểm tra.
\end{enumerate}
Lưu ý: $\overline{A}$ là "$A$ không xảy ra", không phải biến cố "ngược nghĩa" theo cảm nhận. Ví dụ, khi gieo xúc xắc, biến cố đối của "ra số chấm lớn hơn $4$" là "ra số chấm không lớn hơn $4$", không phải "ra số chấm nhỏ hơn $4$".
